\documentclass[12pt]{article}
\usepackage[utf8]{inputenc}
\usepackage[T1]{fontenc}
\usepackage{hyperref}
\title{Usaghade Language Resources for K-12 Curriculum Development}
\author{Compiled by Claude Sonnet 5 (claude-sonnet-5) for Andr\'as Kornai}
\date{}
\begin{document}
\maketitle

\section{Introduction}
Usaghade (ISO 639-3: usk; also Isangele) is a small Lower Cross River language spoken mainly in southwestern Cameroon, with some speakers in Cross River State, Nigeria. Sources give about 10,000 speakers (a dated figure). Very little published material was found in search; this dossier is short on purpose.

\section{Primary Educational Resources}
No primers, textbooks or school materials in Usaghade were found. One source states the language is not taught in schools.

\section{Linguistic and Reference Materials}
\begin{itemize}
\item Wikipedia overview of the Usaghade language: \url{https://en.wikipedia.org/wiki/Usaghade_language}
\item Article on language contact and morphology of Usaghade (University of Bologna journal, \emph{Typology at the Crossroads}): \url{https://typologyatcrossroads.unibo.it/article/view/16309}
\end{itemize}
No dictionary, grammar or orthography was found in the search results.

\section{Religious and Community Resources}
Joshua Project reports that Bible translation has begun but that no complete Scripture exists (mission-oriented source, approximate):
\begin{itemize}
\item Joshua Project, Usakade in Nigeria: \url{https://www.joshuaproject.net/people_groups/print/15742/NI}
\item Joshua Project, people group profile: \url{https://joshuaproject.net/people_groups/15742}
\item Joshua Project, language page: \url{https://joshuaproject.net/index.php/languages/usk}
\end{itemize}

\section{Digital and Audio Resources}
Global Recordings Network has a language page for usk, which notes that no audio recordings are available: \url{https://globalrecordings.net/en/language/usk}

\section{Curriculum Development Context}
The language is endangered and apparently not used in schooling. A curriculum would have to begin with documentation work and orthography development, in cooperation with the speaker community. The contact situation with Bantu languages such as Londo is documented in the Bologna article above.
\end{document}